% Generated draft: requires mathematical and source review.
\documentclass[11pt]{article}
\usepackage[margin=25mm]{geometry}
\usepackage{fontspec}
\setmainfont{Latin Modern Roman}
\usepackage{amsmath,amssymb,mathtools,unicode-math}
\setmathfont{Latin Modern Math}
\usepackage{enumitem,needspace,xurl,hyperref}
\setlength{\parindent}{0pt}
\setlength{\parskip}{5pt}
\setlength{\emergencystretch}{3em}
\newcommand{\N}{\mathbb{N}}
\newcommand{\Z}{\mathbb{Z}}
\newcommand{\Q}{\mathbb{Q}}
\newcommand{\R}{\mathbb{R}}
\newcommand{\C}{\mathbb{C}}
\newcommand{\F}{\mathbb{F}}
\newcommand{\A}{\mathbb{A}}
\newcommand{\PP}{\mathbb{P}}
\newcommand{\E}{\mathbb{E}}
\newcommand{\eps}{\varepsilon}
\newcommand{\dd}{\,\mathrm{d}}
\begin{document}
\section*{UnsolvedMath ID: 8800007}
\textbf{Source problem code:} catalogue entry
\subsection{Definitions and mathematical statement}
Let $K_s$ be a separable closure of the intended field $K$, let $G_K=\operatorname{Gal}(K_s/K)$, and regard $K_s^\times$ as a discrete $G_K$-module with its natural action. The group in the question is continuous Galois cohomology
\[
 H^2(G_K,K_s^\times).
\]
For example, a multiplicative continuous two-cocycle $c$ satisfies
\[
 c(g,h)c(gh,k)=g(c(h,k))c(g,hk),
\]
and cocycles are identified modulo coboundaries
\[
 (\delta b)(g,h)=b(g)\,g(b(h))/b(gh).
\]
The group operation is induced by pointwise multiplication of cocycles.

The supplied problem is to turn this cohomology group into an explicitly computational group. Its usual Brauer-group realization describes classes of finite-dimensional central simple $K$-algebras up to Brauer equivalence, with tensor product as the group law. The research summary specifically invokes number fields and local-invariant/class-field-theoretic descriptions, but the original sentence does not specify the field class or input representation.

A precise computational target would need an effective representation of $K$, finite encodings of cohomology classes, and prescribed operations and tests, such as multiplication, inverse, equality or triviality, and conversion between algebraic, cocycle, and local-invariant data. The source does not require a particular encoding, algorithm, complexity bound, or uniformity in $K$. These possible tasks explain what must be specified; they are not an invented list already solved or a replacement theorem. A theoretical classification through local invariants alone is not silently identified with the requested practical realization.
\subsection{Short English statement}
Develop explicit representations and algorithms for the second Galois cohomology group with multiplicative separable-closure coefficients. The field class and intended operations need specification.

The supplied summary cites the Brauer-group and local-invariant foundations of class field theory, calling them Hasse--Brauer--Noether context, but says no dedicated primary resolution of the practical program was verified. Its OPEN-TRIAGE label describes an unverified research-program request, not an impossibility theorem or a new survey conclusion.
\subsection{Sources}
\begin{itemize}
\item Local UnsolvedMath registry, permanent ID 8800007, source code catalogue entry, ``Computational realization of a second cohomology group,'' proposed by Gerhard Frey. Dataset curation and metadata: Ulam AI and UnsolvedMath Contributors, CC BY 4.0.
\item IPAM, \emph{Workshop on Number Theory and Cryptography --- Open Problems} (2006), Attachment 6219, PDF page 31, Problem 6, schedule locator \url{https://www.ipam.ucla.edu/programs/workshops/workshop-i-number-theory-and-cryptography-open-problems/?tab=schedule}; extraction recorded as pdf-and-legacy-ppt visual audit, accessed 2026-07-29.
\item The supplied summary names Hasse--Brauer--Noether/class field theory without a separate primary reference or exact computational theorem.
\item Source provenance reports inspection of 19 surviving attachments, comprising 17 PDFs and two legacy PPT decks, excluding animation/build repetitions and solved definitions. This is supplied provenance, not an inspection claimed in this draft.
\item Current primary-literature status audit and redistribution terms NEEDS\_REVIEW for the publicly accessible author/publisher copy, separately from the dataset licence.
\end{itemize}
\end{document}
